\documentclass[a4paper]{article}
% \oddsidemargin -0.4in \evensidemargin -0.4in \textwidth 10in
% \topmargin -.65in \textheight 24cm


% \setlength{\textwidth}{7in}

%\usepackage[utf8x]{inputenc}
%\usepackage[T1]{fontenc}
\usepackage{tikz,multicol}
\usepackage{amsmath, amssymb}
% \usepackage{lastpage}
\usepackage{multicol}
% \usepackage{fp}
\usepackage{totcount}
\regtotcounter{AMCquestionaff}

\usepackage[box]{automultiplechoice}
\usepackage{multicol}

\newcommand{\amark}[1]{\scoring{mz=#1}
   (#1 point\ifdim#1pt=1pt%
   \else
      s%
   \fi) 
}

\newcommand{\push}{\hspace*{\stretch{1}}}

%******************** Course Details *********************************
\newcommand{\testname}{Quiz III}
\newcommand{\courseno}{MA 324}
\newcommand{\coursename}{Statistical Inference and Multivariate Analysis}
\newcommand{\instiname}{IIT GUWAHATI}
\newcommand{\testdate}{April 08, 2026}
\newcommand{\testtime}{16:05--16:50 IST}
% \newcommand{\semester}{2025-26 Odd}


\begin{document}

\AMCrandomseed{1237893}

\def\multiSymbole{}
\def\AMCbeginQuestion#1#2{\par\noindent{\bf Question #1:} #2\hspace*{-1mm}}
\def\AMCformQuestion#1{{\bf Question #1:}}    
% 555555
\AMCsetScoreZoneAnswerSheet{width=1.5em,height=1.5ex,depth=.5ex,position=question}

\AMCformHSpace=1em
\AMCformVSpace=5ex

%\setdefaultgroupmode{withoutreplacement}
%\scoring{default}

%GATE 22: q54
\element{mcqgroup}{
   \begin{questionmult}{q1} \amark{2}
      Let a linear model $Y = \beta_0 + \beta_1X + \epsilon$ be fitted to the
      following data, where $\epsilon$ is normally distributed with mean $0$ and
      unknown variance $\sigma^2 > 0$.

      $$
      \begin{array}{|c|c|c|c|c|c|}
         \hline
         x_i & 0 & 1 & 2 & 3 & 4 \\
         \hline
         y_i & 3 & 4 & 5 & 6 & 7 \\
         \hline
      \end{array}
      $$

      Let $\widehat{Y}_0$ denote the ordinary least-square estimator of $Y$ at $X =
      6$. If $E \left( \widehat{Y} \right)=c_0\beta_0+c_1\beta_1$ and $Var\left(
      \widehat{Y}_0 \right) = d\sigma^2$, then which of the
      following is/are true?
      % the value of the
      % real constant $c$ (\textbf{rounded off to one decimal place}) is equal to
      % \underline{\hspace{2cm}}
      \begin{multicols}{2}
         \begin{choices}
            \wrongchoice{$c_1=6,\, c_2=1,\, 1.7\le d\le 1.9$}
            \correctchoice{$c_1=1,\, c_2=6,\, 1.7\le d\le 1.9$}
            \correctchoice{$c_1+c_2+d\le 8.8$}
            \correctchoice{$3.1\le c_2-c_1-d \le 3.5$}
         \end{choices}
      \end{multicols}
   \end{questionmult}
}

%GATE 21: q55
\element{mcqgroup}{
   \begin{questionmult}{q2} \amark{3}
       Let $y_i = \beta_0 + \beta_1x_{1,i} + \beta_2x_{2,i} + \cdots +
       \beta_{22}x_{22,i} + \epsilon_i, i = 1, 2, \cdots, 123$, where,
       $x_{1,i}, \cdots, x_{22,i}$ are fixed input (independent) variables; and,
       for $j=0,1,2, \cdots, 22$;  $\beta_j$ are unknown parameters, and
       $\epsilon_i$'s are independent and identically distributed normal random
       variables with mean zero and finite variance. If $SS_{Reg} = 338.92$ and
       $SS_T = 522.30$ and the value of $R_{adj}^2$ (adjusted $R^2$) equals to
       $t$. Then which of the following statements is/are true?
       %% ANS: 0.57     
       \begin{multicols}{2}
          \begin{choices}
             \wrongchoice{$0.61 < t < 0.69$}
             \wrongchoice{$0.43 < t < 0.54$}
             \correctchoice{$0.51 < t < 0.59$}
             \correctchoice{$0.53 \leq t < 0.63$}
          \end{choices}
       \end{multicols}
    \end{questionmult}
 }

%ISSexam: 2018II: Q42
\element{mcqgroup}{
	\begin{questionmult}{q3} \amark{2}
      Consider a multiple linear regression problem. Which of the following statements is/are true?
      \begin{choices}
         \wrongchoice{A leverage point is always influential point}
         \correctchoice{PRESS is a measure of how well a regression model perform in predicting new observations}
         \correctchoice{Q-Q plot is used to test for normality}
         \correctchoice{A linear model is linear in the parameters, but not necessarily in the independent variables}
      \end{choices}
   \end{questionmult}
}

%Seber-Lee: Page-95, Q-5
\element{mcqgroup}{
   \begin{questionmult}{q4} \amark{3}
      Consider the multiple linear regression problem 
      $$ Y = X \beta + \epsilon, \mbox{ and } \hat{Y} = X \hat{\beta}$$
      where $Y = (Y_1, \cdots, Y_n)'$ is a $n\times 1$ vector of response
      variable, $X$ is a $n \times p$ design matrix with full rank (largest
      possible rank), $\beta$ is a $p \times 1$ vector of parameters, and
      $\epsilon$ is a $n \times 1$ vector of errors; $n>p$. The first element of
      $\beta$ is the intercept of the model. Then which of the following
      statements is/are true?
      \begin{multicols}{2}
         \begin{choices}
            \wrongchoice{$\sum_{i=1}^{n} Var(\hat{Y}_i) = \sigma^2 n$}
            \wrongchoice{$\sum_{i=1}^{n} Var(\hat{Y}_i) < \sigma^2 p$}
            \correctchoice{$\sum_{i=1}^{n} Var(\hat{Y}_i) < \sigma^2 n$}
            \correctchoice{$\sum_{i=1}^{n} Var(\hat{Y}_i) = \sigma^2 p$}
         \end{choices}
      \end{multicols}
   \end{questionmult}
}

%GATE 25: q54
\element{mcqgroup}{
   \begin{questionmult}{q5} \amark{3}
      For $Y \in \mathbb{R}^n, X \in \mathbb{R}^{n \times p}$, and $\beta \in
      \mathbb{R}^p$, consider a regression model
      $$Y = X \beta + \epsilon,$$
      where $\epsilon$ has an $n$-dimensional multivariate normal distribution
      with zero mean vector and identity covariance matrix. Let $I_p$ denote the
      identity matrix of order $p$. For $\lambda > 0$, let
      $$\widehat{\beta}_n = (X^T X + \lambda I_p)^{-1} X^T Y,$$
      be an estimator of $\beta$. Then which of the following options is/are
      correct?
      \begin{choices}
         \wrongchoice{$\widehat{\beta}_n$ is an unbiased estimator of $\beta$}
         \correctchoice{$(X^T X + \lambda I_p)$ is a positive definite matrix}
         \correctchoice{$\widehat{\beta}_n$ has a multivariate normal distribution}
         \wrongchoice{$\text{Var}(\widehat{\beta}_n) = (X^T X + \lambda I_p)^{-1}$}
      \end{choices}
   \end{questionmult}
}

%
\element{mcqgroup}{
   \begin{questionmult}{q6} \amark{2}
      Consider the linear model
      $$ \begin{aligned} 
         Y_1 &= \beta_1 + \beta_2 + \epsilon_1 \\ 
         Y_2 &= \beta_1 + 2\beta_2 + \epsilon_2 \\ 
         Y_3 &= \beta_1 + c\beta_2 + \epsilon_3, 
      \end{aligned} $$
      where $\beta_1, \beta_2 \in \mathbb{R}$ are unknown parameters, and the
      uncorrelated errors $\epsilon_i, i = 1, 2, 3$ have zero mean and finite
      variance $\sigma^2(> 0)$. The constant $c$ is such that $\hat{\beta}_1$
      and $\hat{\beta}_2$ are uncorrelated, where $\hat{\beta}_1$ and
      $\hat{\beta}_2$ are the best linear unbiased estimators of $\beta_1$ and
      $\beta_2$, respectively. Then $(\text{Var}(\hat{\beta}_1),
      \text{Var}(\hat{\beta}_2))$ equals
      \begin{multicols}{4}
         \begin{choices}
            \correctchoice{$\left( \frac{\sigma^2}{3}, \frac{\sigma^2}{14} \right)$}
            \wrongchoice{$(3\sigma^2, \frac{\sigma^2}{14})$}
            \wrongchoice{$\left( \frac{\sigma^2}{14}, \frac{\sigma^2}{3} \right)$}
            \wrongchoice{$(\frac{\sigma^2}{3}, 14\sigma^2)$}
         \end{choices}
      \end{multicols}
   \end{questionmult}
}


% == == == == == == == == == == == == == == == == == == == ==
% This is for defining default grading strategy for questionmult
% \scoringDefaultM{formula=(NMC>0 ? 0 : M*NBC/NB)}
% == == == == == == == == == == == == == == == == == == == ==

\onecopy{85}{

%% beginning of the test sheet header:

\noindent\textsc{{\normalsize\courseno \push \testname}\smallskip\\
{\normalsize\coursename \push \testtime}\smallskip\\
{\normalsize\instiname \push \testdate}}

%\AMCcleardoublepage    

\AMCformBegin

\AMCboxDimensions{shape=oval}

\begin{center}
   \em Answers must be given \textbf{exclusively on this sheet}: answers given on the other sheets
   \textbf{will be ignored}. This question paper has \textbf{\total{AMCquestionaff} questions} and 2
   pages.
\end{center}

\noindent\hrulefill\raisebox{-2pt}{\footnotesize\fbox{\textsc{Name and Roll Number}}}\hrulefill

%%% beginning of the answer sheet header

\noindent
\begin{minipage}{0.5\linewidth}
Please mark your Roll No:\\[2mm]
   \AMCcode{rollno}{9}
\end{minipage}
\begin{minipage}{0.45\linewidth}
   \namefield{\fbox{  
         \begin{minipage}{\linewidth}
         Name of the Student:

         \vspace*{.3cm} \namefielddots

         Roll Number of the Student:

         \vspace*{.3cm}\namefielddots

         Signature of the Student:
            
         \vspace*{.3cm}\namefielddots

         Signature of the Invigilator:
            
         \vspace*{.3cm}\namefielddots
         \end{minipage}
   }}
\end{minipage}

\vspace*{5mm}

\vspace{5mm}

\noindent\hrulefill\raisebox{-2pt}{\footnotesize\fbox{\textsc{Questions and Responses}}}\hrulefill

\vspace*{4mm}

\shufflegroup{mcqgroup}
\insertgroup{mcqgroup}

%\noindent\hrulefill\raisebox{-2pt}{\tiny\fbox{\textsc{Response}}}\hrulefill

%\begin{center}
%\AMCform
%\end{center}

\AMCaddpagesto{2} 

}

\end{document}

